\chapter{SUSTAINABLE DEVELOPMENT GOALS (SDGS)}

The United Nations 2030 Agenda for Sustainable Development provides a shared blueprint for global peace and prosperity, structured around 17 Sustainable Development Goals (SDGs) encompassing social, economic, and environmental targets. Engineering projects and data-driven computational systems bear an ethical responsibility to align technical methodologies with sustainable development objectives. 

Based on the official United Nations SDG targets documented in the institutional reference materials, this case study directly addresses two core Sustainable Development Goals: \textbf{SDG 9 (Industry, Innovation, Technology \& Infrastructure)} and \textbf{SDG 8 (Decent Work and Economic Growth)}.

\section{MAPPING TO SUSTAINABLE DEVELOPMENT GOALS}
Table~\ref{tab:sdg_mapping_full} delineates the explicit relationship between the machine learning deliverables executed in this case study and the corresponding UN SDG targets and institutional criteria.

\begin{table}[H]
\centering
\small
\begin{tabular}{|p{3.2cm}|p{2.4cm}|p{8.5cm}|}
\hline
\textbf{Goal Name} & \textbf{Official Target} & \textbf{Project-Specific Technical Justification} \\
\hline
\textbf{SDG 9: Industry, Innovation, Technology \& Infrastructure} & 
\textbf{Target 9.5:} Enhance scientific research and upgrade industrial technological capabilities. \newline \newline
\textbf{Target 9.c:} Significantly increase access to ICT and affordable internet access. & 
The case study develops an advanced computational machine learning framework that upgrades the analytical and operational capabilities of the telecommunications service industry. By establishing a predictive system capable of identifying root factors behind customer defection, the system aids telecommunications providers in stabilizing digital communication networks, optimizing network resource allocation, and maintaining affordable, high-quality broadband access for subscribers. \\
\hline
\textbf{SDG 8: Decent Work and Economic Growth} & 
\textbf{Target 8.2:} Achieve higher levels of economic productivity through technological upgrading and innovation. & 
Uncontrolled customer attrition causes severe operational inefficiency, leading to wasteful expenditures on redundant customer re-acquisition. By implementing data-driven predictive classification and customer segmentation, the project enhances business productivity, enables optimized workforce resource allocation in customer care departments, and fosters sustainable, resilient commercial operations within the knowledge-based digital services economy. \\
\hline
\end{tabular}
\caption{Comprehensive Project Alignment with United Nations Sustainable Development Goals}
\label{tab:sdg_mapping_full}
\end{table}

\section{DETAILED JUSTIFICATION OF GOAL ALIGNMENT}

\subsection{Alignment with SDG 9: Industry, Innovation and Infrastructure}
Target 9.5 mandates that member states upgrade industrial technical capabilities and support scientific innovation across productive sectors. Telecommunications forms the critical physical and digital backbone of modern digital societies. However, telecommunications infrastructure cannot remain viable if operators suffer from volatile customer attrition that disrupts capital expenditure planning.

By engineering a transparent, leakage-free machine learning architecture that models customer attrition patterns across broadband (DSL, Fiber optic) and telephony networks, this study directly advances Target 9.5. Furthermore, in relation to Target 9.c (expanding affordable internet access), our empirical segmentation uncovered that early-tenure fiber optic customers exhibit an alarming 44.0\% churn rate driven by billing friction. Implementing targeted retention measures (such as transparent onboarding and payment assistance) directly helps maintain subscribers on high-speed internet connections, preventing involuntary digital disenfranchisement.

\subsection{Alignment with SDG 8: Decent Work and Economic Growth}
Target 8.2 emphasizes sustained economic productivity through technological diversification, digitization, and innovation. In conventional customer support operations, customer retention representatives operate under reactive, stressful conditions, attempting to salvage customer relationships only after a cancellation request has been lodged.

The deployment of an automated risk-scoring system fundamentally transforms this workflow into a proactive, structured operation. Customer care personnel can be provided with pre-computed churn probability alerts and behavioral cluster profiles, enabling targeted, constructive, and empathetic customer engagement. This eliminates wasted corporate resources on shotgun marketing campaigns, elevates labor productivity, and stabilizes recurring enterprise revenue in support of sustained economic performance.
